% Shared exercise chunk: l2-threshold-dependence
\begin{exercisechunk}
\item \textbf{Dependent samples and matching marginals.}
\label[exercise]{ex:threshold-dependence}
\exerciseprofile{2}{2}{2}{\exproof\quad\excalculation\quad\exconstruction\quad\exinterpretation}{%
Construct dependent samples that satisfy conditional coverage and show why
matching marginals alone do not ensure improvement with sample size.}
\begin{enumerate}[(a)]
\item Construct a sampling procedure whose training indices are not independent
and verify the conditional-coverage inequality in
\cref{ex:threshold-coverage} for some $\alpha\in(0,1)$.
\item Let $N=2$, $f=f_1$, and $X_1=\cdots=X_m=Z$ for $m\ge1$, where
$Z$ is uniform on $\{1,2\}$. Let the test index be independently uniform.
Compute the learner's risk for every $m\ge1$. Use this expression to determine
whether increasing $m$ reduces the risk in this example.
\end{enumerate}

Explain why the first construction satisfies conditional coverage while the
second has the correct marginal distribution for every training index but
does not.
\end{exercisechunk}
